% =================================================================
% Apresentacao - Validador de Bases de Referencia (Plugin QGIS)
% Compilar com: pdflatex apresentacao.tex (rodar 2x para sumario)
% =================================================================
\documentclass[10pt,aspectratio=169]{beamer}

\usepackage[utf8]{inputenc}
\usepackage[T1]{fontenc}
\usepackage[brazil]{babel}
\usepackage{lmodern}
\usepackage{tikz}
\usepackage{listings}
\usepackage{xcolor}
\usepackage{booktabs}
\usepackage{array}
\usepackage{hyperref}

\usetheme{Madrid}
\usecolortheme{seahorse}
\setbeamertemplate{navigation symbols}{}
\setbeamertemplate{footline}[frame number]
\setbeamercolor{frametitle}{bg=structure.fg!85,fg=white}
\setbeamercolor{title}{bg=structure.fg!85,fg=white}
\setbeamercolor{block title}{bg=structure.fg!80,fg=white}
\setbeamercolor{block body}{bg=structure.fg!8,fg=black}

% --- Codigo ---
\definecolor{codegray}{rgb}{0.95,0.95,0.95}
\definecolor{codekey}{rgb}{0.0,0.0,0.6}
\definecolor{codestr}{rgb}{0.6,0.1,0.1}
\definecolor{codecom}{rgb}{0.2,0.5,0.2}
\lstdefinestyle{py}{
    backgroundcolor=\color{codegray},
    basicstyle=\ttfamily\scriptsize,
    keywordstyle=\color{codekey}\bfseries,
    stringstyle=\color{codestr},
    commentstyle=\color{codecom}\itshape,
    language=Python,
    breaklines=true,
    showstringspaces=false,
    columns=fullflexible,
    frame=single,
    framerule=0pt,
}

% --- TikZ ---
\usetikzlibrary{shapes.geometric,arrows.meta,positioning,fit,backgrounds}
\tikzset{
  bloco/.style={rectangle,rounded corners=3pt,draw=structure.fg,
                fill=structure.fg!10,thick,minimum height=10mm,
                minimum width=22mm,align=center,font=\scriptsize},
  flecha/.style={-Stealth,thick,structure.fg},
  destaque/.style={rectangle,rounded corners=3pt,draw=red!60!black,
                   fill=red!10,thick,minimum height=10mm,align=center,
                   font=\scriptsize\bfseries}
}

% --- Metadados ---
\title[Validador de Bases]{Validador de Bases de Referência\\
\large Plugin QGIS para Validação e Adequação Geoespacial}
\subtitle{Funcionalidades, fluxos e bibliotecas}
\author{Equipe de Desenvolvimento}
\institute{Apresentação técnica para discussão de portabilidade Web}
\date{\today}

\begin{document}

% ============================================================
\begin{frame}
  \titlepage
\end{frame}

% ============================================================
\begin{frame}{Sumário}
  \tableofcontents
\end{frame}

% ============================================================
\section{Visão geral}
% ============================================================
\begin{frame}{O que é o aplicativo}
  \begin{block}{Objetivo}
    Plugin QGIS escrito em Python que oferece um ambiente único para:
    \begin{itemize}
      \item \textbf{Validar a qualidade} de bases geográficas de referência (temática + posicional);
      \item \textbf{Adequar} bases ao padrão da \emph{Análise Dinamizada do CAR};
      \item \textbf{Categorizar a hidrografia} e \textbf{gerar APP hídrica} conforme a Lei n.\ 12.651/2012.
    \end{itemize}
  \end{block}

  \begin{block}{Arquitetura macro}
    \begin{description}
      \item[\texttt{plugin\_main.py}] Entrada do plugin; registra ícone e abre a \texttt{MainWindow}.
      \item[\texttt{ui/}] Interface (\texttt{QMainWindow}, \texttt{QTabWidget}, diálogos, canvases de mapa).
      \item[\texttt{core/}] Algoritmos: amostragem, matriz de confusão, adequação, APP hídrica.
      \item[\texttt{config/}] \texttt{config.json} (serviços XYZ, parâmetros padrão).
    \end{description}
  \end{block}
\end{frame}

% ============================================================
\begin{frame}{Estrutura de abas (após refatoração)}
  \begin{center}
    \begin{tikzpicture}[node distance=8mm and 6mm]
      \node[bloco] (raiz) {MainWindow};
      \node[bloco, below left=of raiz, xshift=-15mm] (aq)  {Avaliação\\da Qualidade};
      \node[bloco, below right=of raiz, xshift=15mm] (ab) {Adequação\\de Bases};

      \node[bloco, below=of aq, xshift=-15mm] (uso) {Uso do Solo};
      \node[bloco, below=of aq, xshift=15mm]  (out) {Outras Bases};

      \node[bloco, below=of ab, xshift=-15mm] (camad)  {Camadas};
      \node[bloco, below=of ab]                (padrao) {Padrão de\\Adequação};
      \node[bloco, below=of ab, xshift=15mm]   (apph)   {Categorizar Hidrog.\\\& Gerar APP};

      \draw[flecha] (raiz) -- (aq);
      \draw[flecha] (raiz) -- (ab);
      \draw[flecha] (aq) -- (uso);
      \draw[flecha] (aq) -- (out);
      \draw[flecha] (ab) -- (camad);
      \draw[flecha] (ab) -- (padrao);
      \draw[flecha] (ab) -- (apph);
    \end{tikzpicture}
  \end{center}
\end{frame}

% ============================================================
\section{Bibliotecas e tecnologias}
% ============================================================
\begin{frame}{Stack tecnológico}
  \begin{columns}[T,onlytextwidth]
    \begin{column}{0.5\textwidth}
      \textbf{Núcleo Python}
      \begin{itemize}
        \item \texttt{Python 3.9+}
        \item \texttt{PyQt5} (Qt 5): UI, sinais/slots, \texttt{QDialog}, \texttt{QTabWidget}, \texttt{QSplitter}.
        \item \texttt{numpy} (cálculos de matriz, índices).
        \item \texttt{requests} (Planet API).
        \item \texttt{earthengine-api} (\texttt{ee}) — opcional, para GEE.
      \end{itemize}
    \end{column}
    \begin{column}{0.5\textwidth}
      \textbf{QGIS API (PyQGIS)}
      \begin{itemize}
        \item \texttt{QgsVectorLayer}, \texttt{QgsRasterLayer}.
        \item \texttt{QgsMapCanvas}, \texttt{QgsRubberBand}, \texttt{QgsMapTool}.
        \item \texttt{QgsVectorFileWriter} (\emph{create}).
        \item \texttt{QgsDistanceArea} (medições geodésicas).
        \item \texttt{QgsCoordinateReferenceSystem} / \texttt{Transform}.
        \item \texttt{processing} (\texttt{native:buffer}, \texttt{difference}, \texttt{dissolve}, \texttt{intersection}, \texttt{subdivide}\ldots).
      \end{itemize}
    \end{column}
  \end{columns}
  \vfill
  \begin{alertblock}{Para a versão Web}
    O \textbf{núcleo geoespacial} pode ser portado para \emph{GeoPandas + Shapely + PyProj + Fiona}, expondo as mesmas regras de negócio sob FastAPI/Flask. As funções já estão isoladas em \texttt{core/} — boa porta de entrada.
  \end{alertblock}
\end{frame}

% ============================================================
\section{Avaliação de Qualidade — Uso do Solo}
% ============================================================
\begin{frame}{Visão da sub-aba ``Uso do Solo''}
  Tudo em uma única tela, com mapa à direita e blocos à esquerda:
  \begin{enumerate}
    \item \textbf{Camadas} — carga do shapefile de uso e cobertura + \emph{mapa de fundo} (Google, Esri, Planet, GEE SPOT/Landsat 2008).
    \item \textbf{Tipo de Amostragem} — escolha da estratégia + parâmetros simplificados (avançados via engrenagem).
    \item \textbf{Calcular pontos} — gera amostra e desenha no mapa em vermelho.
    \item \textbf{Resumo do cálculo} — tabela com totais por classe.
    \item \textbf{Rotulagem} — navegação ponto a ponto comparando classificação vs.\ verdade.
    \item \textbf{Matriz de Confusão} (botão $\rightarrow$ modal).
  \end{enumerate}
\end{frame}

% ============================================================
\begin{frame}{Fluxo de Uso do Solo}
  \begin{center}
    \begin{tikzpicture}[node distance=10mm and 6mm]
      \node[bloco] (a) {Carregar\\uso do solo};
      \node[bloco, right=of a] (b) {Escolher\\campo classe};
      \node[bloco, right=of b] (c) {Estratégia +\\parâmetros};
      \node[bloco, right=of c] (d) {Calcular\\pontos};
      \node[bloco, below=of d] (e) {Rotular ponto\\a ponto};
      \node[bloco, left=of e] (f) {Mapa de fundo\\(Google/Planet/GEE)};
      \node[destaque, left=of f] (g) {Matriz\\de Confusão};
      \draw[flecha] (a)--(b); \draw[flecha](b)--(c); \draw[flecha](c)--(d);
      \draw[flecha] (d)--(e); \draw[flecha](e)--(f); \draw[flecha](f)--(g);
    \end{tikzpicture}
  \end{center}
\end{frame}

% ============================================================
\begin{frame}{Estratégias de amostragem}
  Implementadas em \texttt{core/amostragem.py}.
  \begin{block}{Tamanho da amostra}
    Fórmula multinomial de Congalton \& Green (1999):
    \[
      n = \frac{B \cdot \Pi_i (1-\Pi_i)}{b_i^{\,2}}
    \]
    onde $B$ depende de $\alpha$ (chi-quadrado) e $b_i$ é a precisão desejada.
  \end{block}
  \begin{columns}[T,onlytextwidth]
    \begin{column}{0.5\textwidth}
      \textbf{Distribuição}
      \begin{itemize}
        \item Estratificada \emph{proporcional} (peso pela área da classe).
        \item Estratificada \emph{igualitária}.
        \item Aleatória simples.
        \item Sistemática por grade (espaçamento em metros).
      \end{itemize}
    \end{column}
    \begin{column}{0.5\textwidth}
      \textbf{Detalhes}
      \begin{itemize}
        \item Área pode vir de \emph{campo} existente ou ser calculada via \texttt{QgsDistanceArea} geodésico.
        \item Resultado salvo em \texttt{area\_ha} se ``Calcular Área (ha)''.
        \item Sistemática usa só \emph{bbox} (rápida).
      \end{itemize}
    \end{column}
  \end{columns}
\end{frame}

% ============================================================
\begin{frame}{Mapa de fundo unificado}
  Compartilhado entre Uso do Solo e Outras Bases via \texttt{ui/helpers\_basemap.py}:
  \begin{itemize}
    \item \textbf{Google Satélite} (padrão), \textbf{Esri Satélite} — \texttt{QgsRasterLayer} XYZ via WMS provider.
    \item \textbf{Planet (login)} — diálogo dedicado:
      \begin{itemize}
        \item Login por e-mail/senha (Basic Auth + fallback API key).
        \item Lista \emph{mosaics} filtrados por \texttt{name\_contains="normalized\_analytic"}.
        \item Gera URL de tile $\rightarrow$ camada XYZ.
        \item Lembra credenciais durante a sessão.
      \end{itemize}
    \item \textbf{GEE SPOT 2008 / Landsat 2008} — usa \texttt{earthengine-api}, gera \texttt{getMapId()} e expõe XYZ temporário.
  \end{itemize}
  \begin{exampleblock}{Implicação para o Web}
    Padrão de tiles XYZ é universal — fácil migrar para MapLibre/Leaflet usando os mesmos templates de URL.
  \end{exampleblock}
\end{frame}

% ============================================================
\begin{frame}{Rotulagem assistida}
  Em \texttt{widget\_uso\_solo.py}:
  \begin{itemize}
    \item Zoom inicial no \emph{extent} de todos os pontos sobre a camada de uso.
    \item Ao iniciar a rotulagem: zoom no ponto atual + camada de uso é ocultada (julgamento ``cego'').
    \item Combo com as classes derivadas da própria camada (não há ``tendência'').
    \item Atributos gravados: \texttt{classificacao} (origem) e \texttt{verdade} (rotulado).
    \item Navegação ``Anterior'' / ``Próximo'' / ``Salvar e próximo''.
  \end{itemize}
  \begin{block}{Razões UX}
    A separação entre classificação original (oculta) e verdade (rotulada) reduz viés do operador, alinhada a Olofsson \emph{et al.} (2014).
  \end{block}
\end{frame}

% ============================================================
\begin{frame}{Matriz de Confusão (modal)}
  Em \texttt{ui/dialog\_matriz\_confusao.py}, com \texttt{core/matriz\_confusao.py}.
  \begin{itemize}
    \item Filtra pontos não rotulados (NULL, vazio, ``NULL''/``none''/``nan'').
    \item Calcula:
      \begin{itemize}
        \item \textbf{Global} (Exatidão global) e \textbf{Kappa} (Cohen, 1960).
        \item Acurácia do \emph{produtor} e do \emph{usuário}.
        \item Erros de \emph{omissão} e \emph{comissão} por classe.
      \end{itemize}
    \item Visual com cores dinâmicas conforme Landis \& Koch (1977).
    \item Exporta CSV.
  \end{itemize}
  \begin{block}{Kappa}
    \[
      \kappa = \frac{p_o - p_e}{1 - p_e}, \quad
      p_o = \tfrac{1}{N}\sum_i n_{ii}, \quad
      p_e = \tfrac{1}{N^2}\sum_i n_{i+}\, n_{+i}
    \]
  \end{block}
\end{frame}

% ============================================================
\section{Avaliação de Qualidade — Outras Bases}
% ============================================================
\begin{frame}{Sub-aba ``Outras Bases''}
  Em \texttt{widget\_quadrantes.py}. Mesmo padrão de layout (mapa à direita, blocos à esquerda) e três grupos:
  \begin{enumerate}
    \item \textbf{Camadas} — base a avaliar + mapa de fundo (Google/Esri/Planet, mesmo helper da outra aba).
    \item \textbf{Amostragem} — tamanho do quadrante (m), nº a sortear, \emph{Gerar quadrantes}.
    \item \textbf{Julgamento dos quadrantes} —
      \begin{itemize}
        \item PEC: medição efêmera (clique-arraste-solta) com \emph{label flutuante} em metros (\texttt{QgsDistanceArea}).
        \item Erros temáticos: spins de omissão / comissão.
        \item Aprovar / Reprovar.
      \end{itemize}
    \item Botão \textbf{Exibir resultados} — abre modal interativo.
  \end{enumerate}
\end{frame}

% ============================================================
\begin{frame}{Geração e sorteio de quadrantes}
  \begin{itemize}
    \item Em \texttt{core/amostragem.py::gerar\_quadrantes}.
    \item Cria malha por \emph{bbox} da camada, usando o tamanho em metros (CRS métrico).
    \item Usa \texttt{QgsSpatialIndex} para filtrar quadrantes que \emph{tocam} a camada avaliada.
    \item Sorteio com \emph{reservoir sampling} (memória O(n) sem armazenar tudo).
    \item Estilo de saída: somente contorno vermelho (não obstrui a imagem).
  \end{itemize}
  \begin{block}{Por que ``só os sorteados''}
    Antes mostrávamos toda a malha (centenas/milhares de quadrantes), o que cobria a tela. Agora a camada de saída já chega filtrada.
  \end{block}
\end{frame}

% ============================================================
\begin{frame}{PEC + Resultados configuráveis}
  Em \texttt{ui/dialog\_resultados\_quadrantes.py}:
  \begin{itemize}
    \item Não há ``parecer'' fixo — usuário escolhe \textbf{escala}, \textbf{classe PEC} (A/B/C), \textbf{LQA} e \textbf{máx.~omissão+comissão} por quadrante; tudo recalcula ao vivo.
    \item Tabela detalhada por quadrante: \emph{ID, julgamento, distância, dentro PEC?, omissão, comissão, soma, dentro O+C?}.
    \item Resumo: distância média/máxima, $\Sigma$ omissão/comissão, LQA observado.
    \item Cores semânticas: verde (dentro), vermelho (fora).
    \item Exporta CSV (com critérios usados no cabeçalho).
  \end{itemize}
  \begin{exampleblock}{Tabela PEC (Lei)}
    Para escala 1:50.000: A=25\,m, B=40\,m, C=50\,m. Outras escalas configuradas em \texttt{PEC\_TABELA}.
  \end{exampleblock}
\end{frame}

% ============================================================
\section{Adequação de Bases}
% ============================================================
\begin{frame}{Estrutura da aba ``Adequação de Bases''}
  Aba única (sem sub-abas), mapa à direita, painel à esquerda dividido em 3 grupos.
  \begin{block}{Grupo Camadas}
    Base a adequar + mapa de fundo (Google [padrão], Esri, Planet).
  \end{block}
  \begin{block}{Grupo Padrão de Adequação}
    \textbf{Análise Dinamizada do CAR} (único padrão atual) + pasta + formato (SHP/GPKG) + \textbf{Tipo de base} + opção \emph{Quebrar polígonos com mais de N vértices} (Dice).
  \end{block}
  \begin{block}{Grupo Categorizar Hidrografia e Gerar APP}
    Botão único que abre um diálogo modal dedicado.
  \end{block}
\end{frame}

% ============================================================
\begin{frame}{Tipos de base e UI dinâmica}
  Página dinâmica via \texttt{QStackedWidget}:
  \begin{itemize}
    \item \textbf{Individuais} (Uso do Solo, Fitofisionomia): gera \emph{N} shapefiles, um por categoria padronizada (\emph{ANTROPIZADO, CONSOLIDADO, VEGETACAO\_ATUAL, VEGETACAO\_2008}; \emph{AML\_FLORESTA/CERRADO/CAMPO}).
    \item \textbf{Unificadas} (APP\_ESPECIAL, RELEVO, SERVIDAO, USO\_RESTRITO): gera \textbf{1} shapefile com coluna \texttt{CLASSE} (1..N).
  \end{itemize}
  \textbf{Fluxo único:}
  \begin{enumerate}
    \item Escolhe o campo classe da camada original.
    \item ``Detectar valores únicos''.
    \item Combo \emph{Categoria a configurar} + lista única de valores (Ctrl/Shift).
    \item Troca de categoria preserva seleções já feitas.
    \item ``Separar e salvar'' / ``Gerar base unificada''.
  \end{enumerate}
\end{frame}

% ============================================================
\begin{frame}[fragile]{Subdivide (equivalente ao Dice)}
  Em \texttt{core/adequacao\_bases.py::subdividir\_poligonos\_complexos}:
  \begin{itemize}
    \item Usa \texttt{native:subdivide} (algoritmo nativo do QGIS).
    \item Parâmetro \texttt{MAX\_NODES} (vértices) — padrão da UI: 500.
    \item Acionado por checkbox no \textbf{Padrão de Adequação}.
    \item Aplicado em memória \emph{antes} da gravação, mantendo \texttt{CLASSE} e demais atributos.
  \end{itemize}
  \begin{lstlisting}[style=py]
out = processing.run("native:subdivide", {
    "INPUT": mem,
    "MAX_NODES": int(max_vertices),
    "OUTPUT": "memory:",
})["OUTPUT"]
  \end{lstlisting}
\end{frame}

% ============================================================
\begin{frame}{Salvamento robusto}
  Detalhes importantes da camada de I/O:
  \begin{itemize}
    \item \texttt{QgsVectorFileWriter.create(\ldots)} com \texttt{wkbType} e \texttt{crs} herdados — evita o conflito clássico \emph{single} vs.\ \emph{multi} (erro código 8).
    \item \texttt{os.path.normpath} para evitar mistura de \texttt{/} e \texttt{\textbackslash{}} no Windows.
    \item Helper \texttt{\_nome\_tipo\_geom(wkb)} usa \texttt{QgsWkbTypes.flatType} para remover sufixos Z/M sem quebrar ``MultiPolygon''.
    \item Sem reprojeção forçada nem ``fix geometries'' (eram lentos e mascarava bugs).
  \end{itemize}
\end{frame}

% ============================================================
\begin{frame}{Categorizar Hidrografia e Gerar APP}
  Diálogo modal \texttt{ui/dialog\_app\_hidrica.py} + \texttt{core/app\_hidrica.py}.
  \begin{itemize}
    \item Aceita 3 camadas (qualquer combinação): \emph{Trecho de Drenagem}, \emph{Trecho de Massa d'Água}, \emph{Massa d'Água}.
    \item Campo categoria (massa) $\rightarrow$ usuário marca quais valores são \textbf{Reservatório Artificial} (CLASSE 7); restantes viram \textbf{Lago/Lagoa Natural} (CLASSE 6).
    \item Larguras conforme Lei n.\ 12.651/2012 — codificadas em \texttt{\_largura\_app\_para\_rio()}.
    \item Saídas \textbf{duplas}, com apenas \texttt{CATEGORIA} + \texttt{CLASSE}:
      \begin{itemize}
        \item \texttt{HIDROGRAFIA\_AAAAMMDD\_hhmm.shp}
        \item \texttt{APP\_HIDRICA\_AAAAMMDD\_hhmm.shp}
      \end{itemize}
  \end{itemize}
\end{frame}

% ============================================================
\begin{frame}{Pipeline de APP Hídrica}
  \begin{center}
    \begin{tikzpicture}[node distance=7mm and 6mm,every node/.style={font=\scriptsize}]
      \node[bloco] (td) {Trecho\\drenagem};
      \node[bloco, right=of td] (tmd) {Trecho de\\massa d'água};
      \node[bloco, right=of tmd] (md) {Massa\\d'água};

      \node[bloco, below=of td]  (b1) {Buffer 30\,m\\(\texttt{native:buffer})};
      \node[bloco, below=of tmd] (b2) {Largura 4A/P\\$\rightarrow$ buffer\\(30/50/100/200/500)};
      \node[bloco, below=of md]  (b3) {Buffer\\30/50/100\,m\\por categoria};

      \node[bloco, below=of b2] (jun) {Junção c/ prioridade\\(\texttt{difference}+\texttt{merge})};

      \node[bloco, below=of jun, xshift=-20mm] (hid) {HIDROGRAFIA\\(\texttt{intersection}\\\& buffer 0.5\,m)};
      \node[destaque, below=of jun, xshift=20mm] (app) {APP final\\(APP $-$ HIDRO)};

      \draw[flecha](td)--(b1); \draw[flecha](tmd)--(b2); \draw[flecha](md)--(b3);
      \draw[flecha](b1)--(jun); \draw[flecha](b2)--(jun); \draw[flecha](b3)--(jun);
      \draw[flecha](jun)--(hid); \draw[flecha](jun)--(app);
    \end{tikzpicture}
  \end{center}
\end{frame}

% ============================================================
\begin{frame}{APP: detalhes finais}
  \begin{itemize}
    \item \textbf{Dissolve por (CLASSE,CATEGORIA) + MultiToSingle} no fim da pipeline:
      \begin{itemize}
        \item Une polígonos contíguos da mesma classe.
        \item Mantém polígonos desconectados como \emph{features} separadas.
      \end{itemize}
    \item CLASSE 1 (trecho $<$ 10\,m) é \emph{atribuída} ao buffer 0.5\,m antes do merge final, evitando ``CLASSE zerada''.
    \item APP final é recortada (\emph{difference}) pela HIDROGRAFIA — APP não sobrepõe o corpo d'água.
    \item Categorias derivadas pela tabela \texttt{CATEGORIAS\_POR\_CLASSE} quando faltam no atributo original.
  \end{itemize}
\end{frame}

% ============================================================
\section{Pensando na versão Web}
% ============================================================
\begin{frame}{Recomendações de portabilidade}
  \begin{columns}[T,onlytextwidth]
    \begin{column}{0.5\textwidth}
      \textbf{Backend}
      \begin{itemize}
        \item \texttt{FastAPI} + \texttt{GeoPandas/Shapely/PyProj/Fiona}.
        \item Mesmos \emph{cores} (\texttt{core/*}) podem ser refatorados para não depender de PyQGIS — separar lógica das chamadas \texttt{processing.run}.
        \item Processamento pesado: filas (\texttt{Celery}/\texttt{rq}) + workers.
        \item Storage: \texttt{S3/MinIO} para shapefiles e GeoPackages.
        \item PostGIS como banco geográfico (opcional).
      \end{itemize}
    \end{column}
    \begin{column}{0.5\textwidth}
      \textbf{Frontend}
      \begin{itemize}
        \item Mapa: \texttt{MapLibre} ou \texttt{Leaflet} — usa os mesmos templates XYZ.
        \item Componentes: \texttt{React}/\texttt{Vue}; estado: \texttt{Zustand}/\texttt{Pinia}.
        \item Drag para medir distância: implementação direta com \emph{turf.js}.
        \item Planet/GEE: backend continua emitindo URLs assinadas; frontend só consome.
      \end{itemize}
    \end{column}
  \end{columns}
  \vfill
  \begin{alertblock}{Reaproveitamento direto}
    Lógica de matriz de confusão, amostragem multinomial, PEC, dissolves/buffers e categorização da hidro \textbf{já} estão separadas da UI — primeiro alvo de migração.
  \end{alertblock}
\end{frame}

% ============================================================
\begin{frame}{Resumo executivo}
  \begin{itemize}
    \item \textbf{Plugin QGIS} maduro, dois grandes blocos: \emph{Avaliação de Qualidade} (uso do solo + outras bases) e \emph{Adequação de Bases} (organização + APP hídrica).
    \item \textbf{Padrões abertos}: XYZ tiles, shapefile/GPKG, CSV, SIRGAS 2000.
    \item \textbf{Algoritmos consolidados}: Congalton \& Green, Cohen Kappa, Landis \& Koch, PEC IBGE, Lei n.\ 12.651/2012.
    \item \textbf{UX consistente}: layout splitter (ferramentas à esquerda, mapa à direita), modais para resultados, helpers compartilhados.
    \item \textbf{Caminho web}: backend Python isolado em \texttt{core/} já facilita a transição.
  \end{itemize}
\end{frame}

% ============================================================
\begin{frame}
  \centering
  {\Huge Obrigado!}\\[6mm]
  Dúvidas e contribuições são bem-vindas.\\[4mm]
  \small Código-fonte: \texttt{validador\_bases\_referencia/}
\end{frame}

\end{document}
